\chapter{Cost Models and Critical Paths}
\label{ch:cost}

The first surprising invoice is a rite of passage. An agent that cost pennies in testing bills four figures in its first production week, and the postmortem always finds the same thing. Nobody could say where the money went. Not which feature, which user, or which step. Just a total.

That is a measurement failure before it is an efficiency failure, and this chapter is about fixing the measurement first.

We start with the token, work out how token usage compounds across prompts, retrieval, and multi-step execution, then decompose the total into components that respond to different fixes. After that we look at latency, which follows different rules. Cost adds up across everything the agent does, while latency adds up only along the critical path. Confusing the two produces optimizations that make an agent cheaper and slower, or faster and ruinous.

One feature of this domain has no analogue in ordinary capacity planning. The resource gets dramatically cheaper while you are planning for it. The price of a fixed level of capability has fallen roughly $1{,}000\times$ over three years \cite{llmflation2024}. Every cost model here is therefore time-indexed, and every optimization has to justify itself against the alternative of waiting six months and paying less for the same thing. Some of them do not survive that test, and this chapter says which.

\section{The Token Economy}

\subsection{The Puzzle}

Consider a concrete budgeting problem. You have a budget of \$1.00 to answer a single question. The language model charges \$2.50 per million input tokens and \$10.00 per million output tokens. Your system prompt is 500 tokens. The user's question is 50 tokens. You intend to retrieve documents to assist the model. Each retrieved document averages 800 tokens. The model’s response will be approximately 300 tokens long.

The question is simple. How many documents can you retrieve before exceeding your budget?

At first glance, this is straightforward arithmetic. But working through the arithmetic exposes the structure of token-based cost and highlights the tradeoffs inherent in agent design. Input tokens and output tokens are priced differently. Input tokens accumulate across system prompts, user messages, retrieved context, and conversation history. Output length is partially under the designer’s control. Every design decision has an associated cost consequence.

We begin by identifying fixed and variable components. Fixed input costs consist of the system prompt (500 tokens) and the user question (50 tokens), for a total of 550 input tokens. Output cost is fixed at 300 tokens. The variable cost arises from document retrieval, since each document contributes 800 input tokens.

Let $d$ denote the number of retrieved documents. Total input tokens are
\[
550 + 800d.
\]

The cost expression is therefore
\[
\text{Cost} = \frac{550 + 800d}{1{,}000{,}000} \times 2.50
             + \frac{300}{1{,}000{,}000} \times 10.00.
\]

Simplifying,
\[
\text{Cost} = 0.001375 + 0.002d + 0.003
            = 0.004375 + 0.002d.
\]

Imposing the budget constraint,
\[
0.004375 + 0.002d \leq 1.00,
\]
which yields
\[
d \leq \frac{0.995625}{0.002} \approx 497.
\]

So the budget permits 497 documents. Retrieving them would be a mistake.

Long contexts degrade retrieval accuracy in position-dependent ways, add latency proportional to length, and make it much harder to say which document an answer came from. The calculation tells you the capacity of the budget and nothing about the quality of spending it. Nearly every question in this chapter has that shape. Affordability and effectiveness are separate axes, and the interesting designs sit where they disagree.

\subsection{Token Pricing}

Language model APIs price computation in terms of tokens. A token corresponds roughly to four characters of English text, or about three-quarters of a word. For example, the sentence “The quick brown fox jumps over the lazy dog” contains nine words and approximately eleven tokens.

Pricing distinguishes between input tokens and output tokens. Input tokens are those sent to the model, including system prompts, user messages, retrieved documents, and conversation history. Output tokens are those generated by the model. Output tokens typically cost two to four times more than input tokens. This asymmetry reflects the underlying computation. Input processing can be parallelized during the prefill phase, whereas output generation is autoregressive and inherently sequential.

Representative prices from late 2025 appear below. Treat the absolute numbers as a snapshot that is already drifting; the spread between rows is the durable part, and it is the part the arguments in this book rely on.

\begin{center}.
\begin{tabular}{lccc}
\toprule
Model & Input (\$/M) & Output (\$/M) & Context \\
\midrule
GPT-4o & \$2.50 & \$10.00 & 128K \\
GPT-4o-mini & \$0.15 & \$0.60 & 128K \\
Claude Sonnet 4 & \$3.00 & \$15.00 & 200K \\
Claude Haiku 4 & \$0.80 & \$4.00 & 200K \\
Gemini 1.5 Pro & \$1.25 & \$5.00 & 1M \\
Gemini Flash-Lite & \$0.075 & \$0.30 & 128K \\
Llama 4 Scout (Groq) & \$0.11 & \$0.34 & 128K \\
DeepSeek-V3 & \$0.27 & \$1.10 & 128K \\
\bottomrule
\end{tabular}
\end{center}

Look at the spread. It is more than a $30\times$ range in input price across rows that are all plausible choices for the same task. Choosing a model is a cost decision of the same magnitude as choosing whether to do the work at all, which is why Chapter~\ref{ch:action-selection} treats model selection as an action the policy makes rather than a configuration constant.

One cost appears in no pricing table. It is the encoding of the messages themselves. Tool schemas and tool results are structured data, and JSON, the default, was designed for machine interchange rather than token thrift. A controlled comparison of token-optimized formats across agentic workloads found meaningful reductions in the tokens spent purely on syntax, on traffic that carries no additional information \cite{notation2026}. For an agent whose context is mostly tool schemas and tool output, that is free money, and it requires no change to the policy.

\subsection{LLMflation: The Deflation of Intelligence}

A central economic property of large language models is rapid cost deflation. Historical analysis shows that for a fixed level of capability, measured by benchmarks such as MMLU, inference costs decrease by approximately an order of magnitude per year \cite{llmflation2024, epochpricing2025}.

When GPT-3 became publicly accessible in 2021, it achieved an MMLU score of approximately 42 at a cost of \$60 per million tokens. By 2024, equivalent performance could be obtained for roughly \$0.06 per million tokens. For higher capability tiers, price declines have been similarly dramatic, with reductions of 60$\times$ or more over 18-month periods.

This deflation arises from several independent factors.

\textbf{Hardware improvements}. Each generation of accelerators delivers more compute per dollar, driven by architectural improvements rather than raw transistor scaling alone.

\textbf{Quantization}. Inference precision has dropped from 16-bit to 8-bit and now to 4-bit for many workloads, reducing memory bandwidth and increasing throughput with minimal quality loss.

\textbf{Software optimization}. Techniques such as continuous batching, Flash Attention, speculative decoding, and PagedAttention substantially improve utilization of existing hardware.

\textbf{Model efficiency}. Models trained on larger and more diverse corpora achieve comparable performance with fewer parameters. A billion-parameter model in 2025 often exceeds the performance of a 175-billion-parameter model from 2022.

\textbf{Competitive pressure}. Open models and multiple hosting providers compress margins and reduce prices across the ecosystem.

The implications for agent design are practical.
\begin{itemize}
\item Cost models must be time-indexed rather than static.
\item Complex optimization infrastructure may have diminishing returns.
\item Previously uneconomical applications can become viable rapidly.
\item Systems should be designed to adopt cheaper models without architectural change.
\end{itemize}

\section{Cost Decomposition}

Agent costs can be decomposed into components with distinct behaviors and optimization strategies.

\subsection{Token Costs}

Token costs dominate most agent workloads
\[
C_{\text{tokens}} =
\sum_{i=1}^{n}
\left(
T_{\text{in}}^{(i)} p_{\text{in}}
+
T_{\text{out}}^{(i)} p_{\text{out}}
\right),
\]
where $n$ is the number of LLM calls.

In multi-turn interactions, input tokens grow with history
\[
T_{\text{in}}^{(i)} =
T_{\text{system}}
+
\sum_{j=1}^{i-1}
\left(
T_{\text{user}}^{(j)} + T_{\text{assistant}}^{(j)}
\right)
+
T_{\text{user}}^{(i)}.
\]

This cumulative replay leads to quadratic growth in total token usage. For example, a 20-turn conversation with 500 tokens per turn results in approximately 100,000 input tokens across all calls.

\subsection{Tool Costs}

Tool invocations incur their own costs
\[
C_{\text{tools}} = \sum_{j=1}^{m} c_{\text{tool}}^{(j)}.
\]

These costs vary widely. Web search APIs may charge cents per query. Database queries are often negligible. Code execution costs depend on runtime. External APIs may dominate overall cost if invoked frequently.

\subsection{Compute Costs}

For self-hosted models
\[
C_{\text{compute}} = T_{\text{GPU}} \cdot p_{\text{GPU}}.
\]

Self-hosting becomes economical only above some volume threshold, and the threshold is workload-specific rather than universal. It depends on utilization, the hardware you can actually keep busy, and the engineering time the deployment consumes. Published total-cost-of-ownership analyses work the comparison in full for particular settings \cite{tcoanalysis2024}; the general lesson is that below the crossover, managed APIs usually win once operational overhead is counted, and that most teams estimate their own utilization optimistically.

\subsection{Total Cost and Unit Cost}

Total cost aggregates components
\[
C_{\text{total}} = C_{\text{tokens}} + C_{\text{tools}} + C_{\text{compute}}.
\]

More informative is unit cost
\[
C_{\text{unit}} = \frac{C_{\text{total}}}{N_{\text{success}}}.
\]

Unit cost captures both efficiency and effectiveness. Systems that fail frequently may appear cheap per request but expensive per successful outcome.

\section{Critical Path Analysis}

Cost alone does not determine usability. Latency determines whether users tolerate an agent.

\subsection{Latency Decomposition}

Total latency decomposes as
\[
T_{\text{total}} = T_{\text{queue}} + T_{\text{critical}}.
\]

For a single model call
\[
T_{\text{LLM}} = T_{\text{TTFT}} + n_{\text{out}} \cdot T_{\text{per-token}}.
\]

\subsection{Critical Path in Multi-Step Agents}

In a $k$-step ReAct agent
\[
T_{\text{critical}} =
\sum_{i=1}^{k}
\left(
T_{\text{LLM}}^{(i)} + T_{\text{tool}}^{(i)}
\right).
\]

Sequential dependencies dominate latency. Parallelization shortens the critical path but may increase total cost.

\subsection{Critical Path in Multi-Agent Systems}

Multi-agent systems trade coordination overhead for parallelism
\[
T_{\text{parallel}} =
\max_i T_{\text{agent}_i} + T_{\text{aggregation}}.
\]

The Anthropic research system \cite{anthropicresearch2025} exploits this structure to reduce latency for complex queries.

\subsection{Latency-Cost Tradeoffs}

Some optimizations improve both latency and cost (caching). Others trade one for the other (batching). Production systems typically optimize for P95 latency rather than mean latency.

\section{Caching Strategies}

Caching is the one optimization that improves cost and latency at the same time.
Everything else in this chapter trades one against the other. That makes caching
the first thing to reach for and, because it introduces staleness and correctness
questions, the thing most likely to be reached for carelessly.

Four layers of reuse exist in an agent stack, and they exploit different
redundancies. Understanding which redundancy each one captures tells you which
will actually fire on your traffic.

\subsection{KV Cache}

During autoregressive generation, every new token attends to every previous
token. Recomputing those key and value projections at each step would make
generation quadratic in sequence length. The KV cache stores them, reducing
generation to linear time at the cost of memory that grows with context length.

For agents this shifts the binding constraint. Compute stops being the limit and
memory bandwidth becomes it, because the cache for a long-context agent can
exceed the model weights themselves. PagedAttention \cite{vllm2023} addresses the
resulting fragmentation by managing the cache in fixed-size blocks, the way an
operating system manages pages, which raises achievable batch sizes substantially.

Agentic workloads stress this differently from chat. They run long contexts,
multi-turn structure, and repeated tool-call formatting, and a recent survey of
system-aware KV optimization catalogs the techniques that follow from those
properties \cite{kvsurvey2026}. Quantization aimed specifically at agentic
inference pushes cache precision lower than chat workloads tolerate, exploiting
the fact that much of an agent's context is structured, repetitive tool traffic
rather than prose \cite{triaxialkv2026}.

One consequence deserves flagging here because it reappears in
Chapter~\ref{ch:speculation}. The KV cache is state, it survives operations the
application regards as cleanup, and an agent that abandons a reasoning branch
does not thereby remove that branch from the cache \cite{abortedkv2026}. Caching
is a performance mechanism with correctness consequences.

\subsection{Prompt Caching}

Prompt caching reuses the prefill computation for a shared prefix across
requests. An agent with a 3{,}000-token system prompt and a stable tool schema
pays for that prefix on every call; caching it means paying full price once and a
reduced rate thereafter.

Providers price this in a characteristic shape. A write costs somewhat more than an uncached token, a read costs dramatically less, and the cache expires after a
short idle period. Savings on cached tokens typically land in the 50--90\% range.

The shape of that pricing dictates prompt layout. Put everything stable at the front, meaning system instructions, tool definitions, and few-shot examples. Put everything variable at the back. A single user-specific token near the beginning
invalidates the entire suffix, so the common mistake of interpolating a user ID
or timestamp into the system prompt can eliminate the cache benefit while
appearing to do nothing.

Let $f$ be the fraction of tokens in the cacheable prefix, $h$ the hit rate,
$r$ the read multiplier, and $w$ the write multiplier. Relative cost is
\[
\frac{C_{\text{cached}}}{C_{\text{base}}}
= (1-f) + f\big[h \cdot r + (1-h)\cdot w\big].
\]
With $r = 0.1$, $w = 1.25$, $f = 0.8$, and $h = 0.9$, this gives roughly $0.30$:
a 70\% reduction. Setting the bracket equal to 1 gives the break-even hit rate,
\[
h^{*} = \frac{w-1}{w-r},
\]
which for these constants is about $0.22$. Below a 22\% hit rate, caching costs
more than it saves. That threshold is low enough that caching is almost always
correct, and high enough that a badly ordered prompt can genuinely lose money.

\subsection{Semantic Caching}

Semantic caching goes further. Embed the incoming query, search for a previously answered query within some similarity threshold, and return the stored response.
A hit skips inference entirely, so the saving is total rather than partial.

Caching designed for chatbots transfers poorly to agents, because an agent's output depends on external data and environmental context that a query embedding does not capture. Caching at the level of the agent's \emph{plan} rather than its response addresses this more directly \cite{plancache2025}.

The threshold is the entire design. Set it loose and the cache answers ``what is
our refund policy for enterprise customers'' with the response for ``what is our
refund policy,'' which is a correctness failure that looks like a latency
improvement. Set it tight and the hit rate collapses. Unlike the lower layers,
semantic caching can return a wrong answer, so it belongs behind the same kind of
gate as any other action that commits to a result, and it should be scoped by
tenant and by permission. A cache shared across users is an access-control
boundary that nobody reviewed.

Semantic caching suits high-volume, low-variance traffic such as support FAQs, documentation lookup, and repeated internal queries. It suits open-ended research
tasks poorly, because the premise of a research agent is that the question has
not been asked before.

\subsection{Semantic Cache Keys Are Fuzzy Hashes}

The threshold problem above has a sharper statement.

Treat a semantic cache key as a fuzzy hash and a tension becomes formal. Cache
hit rates require \emph{locality}, meaning nearby inputs map to nearby keys.
Collision resistance requires the opposite, the avalanche property, meaning a
small input change produces a large key change. A semantic cache cannot have
both, and the trade is inherent rather than an artifact of any particular
embedding \cite{cachecollision2026}.

The security consequence is a key collision attack. An adversary who can shape
queries can engineer a collision and cause the cache to serve one user's
response to another, which is a confidentiality failure wearing the costume of a
performance optimization.

Two design rules follow. Scope semantic caches by tenant and permission so that a
collision cannot cross a trust boundary. And treat a cache hit as a proposal that
passes a gate rather than as an answer, exactly as Chapter~\ref{ch:gates} treats
any other action with a nonzero error rate.

\subsection{Multi-Layer Caching Architecture}

Production systems layer all of these, cheapest check first:

{\footnotesize{\footnotesize
\begin{verbatim}
  query
    |
    v
  [semantic cache]  hit -> return stored response      (100% saved)
    | miss
    v
  [prompt cache]    hit -> reuse prefix prefill        (50-90% on prefix)
    | miss
    v
  [KV cache]        within-request reuse               (quadratic -> linear)
    | miss
    v
  [inference]       full cost
\end{verbatim}
}}

Each layer captures a different redundancy: semantic caching exploits repeated
\emph{questions}, prompt caching exploits repeated \emph{context}, and the KV
cache exploits repeated \emph{attention} inside one generation. A workload with
no repeated questions still benefits from the lower two, which is why the layers
are worth keeping separate rather than collapsing into one hit-rate number.

Multi-agent systems complicate this. When several agents share a large context,
naive cache reuse across candidates can change results rather than merely
accelerate them, since the interaction between candidates carries information
that reuse discards \cite{kvreuse2026}. Measure quality, not just hit rate.

\subsection{The Scheduler Does Not Know You Are an Agent}

One assumption runs underneath all of the above, and it is wrong on shared
infrastructure.

GPU schedulers treat each model call as an independent request. An agent
executing a hundred chained calls therefore has its intermediate state discarded
between steps, so the KV cache built during step 7 is gone by step 8 and gets
rebuilt from scratch. Measured end to end, this request-level abstraction
inflates agent latency by a reported factor of three to eight
\cite{saga2026}.

The fix is to make the workflow rather than the call the schedulable unit.
Capturing the agent's execution graph lets the scheduler predict which KV state
will be reused across tool-call boundaries and keep it resident. This is
program-level scheduling, and it recovers latency that no amount of
prompt-side optimization can reach, because the loss happens below the API.

Related work treats agent networks as a queueing-control problem in which
contextual memory state, including prefix blocks and verified tool outputs,
alters service rates rather than sitting passively as a cache
\cite{backpressure2026}. The framing matters for capacity planning. An agent
fleet's throughput depends on state locality, so scheduling decisions and cache
decisions are the same decision.

If you self-host, this is actionable today. If you use a provider, it explains a
gap between your measured latency and the vendor's published per-call numbers.

\section{Cost Optimization Strategies}

With measurement in place, the interventions below are roughly ordered by return
on effort. The multipliers compose, which is why a system that applies four of
them lands an order of magnitude below the naive implementation.

\textbf{Cascading.} Attempt the task with a small model, escalate on low
confidence. Chapter~\ref{ch:action-selection} derives the escalation threshold
properly; the headline is that cascading exploits the fact that most queries are
easy and pricing is steeply tiered. BudgetMLAgent reports 94\% cost reduction
through multi-model orchestration on its workload \cite{budgetmlagent2024}. The
catch, developed in Chapter~\ref{ch:action-selection}, is that the escalation
decision needs \emph{calibrated} confidence, and most deployed routers do not
have it \cite{ucci2026}.

\textbf{Context management.} Input tokens dominate multi-turn agents because
history is replayed on every call. Trajectory compression cuts what is replayed
\cite{agentdiet2024}; observation masking drops the bulky middles of tool outputs
that the model never referenced. Chapter~\ref{ch:memory} covers the mechanics and
the failure modes, of which there are more than the cost literature suggests.

\textbf{Output control.} Output tokens cost several times input tokens, and
output length is one of the few quantities fully under your control. Asking for
structured output rather than prose, capping reasoning length by task difficulty,
and forbidding restatement of the question are unglamorous and reliably effective.

\textbf{Early stopping.} Agents routinely continue past the point of diminishing
return, taking additional steps that do not change the answer. The gain here
comes from deciding \emph{when to stop} rather than from adding capability
\cite{budgettooluse2026}, and Chapter~\ref{ch:budgets} makes stopping a policy
decision with an explicit criterion.

\textbf{Batching.} Amortizes fixed overhead across requests at the cost of
latency for the earliest arrival. Correct for offline and background work,
usually wrong for interactive paths.

\textbf{Encoding.} Switching tool-call payloads away from verbose JSON reduces
tokens spent on syntax \cite{notation2026}. Small percentage, zero risk, applies
to every call.

\textbf{Topology.} In multi-agent systems the communication graph is itself a
cost variable. Choosing backbone models per role and then optimizing the topology
under a budget \cite{scmas2026}, or routing dynamically on observed progress
\cite{progrouter2026}, addresses a cost that single-agent thinking cannot see.

\section{Cost Monitoring and Attribution}

An unattributed total is not actionable. If all you know is that last month cost
\$40{,}000, every optimization proposal is a guess. Attribution converts the
number into a decision.

Instrument along four dimensions, all of which should be available on every trace
event described in Chapter~\ref{ch:traces}.

\begin{center}
\begin{tabular}{lll}
\toprule
Dimension & Question it answers & Typical finding \\
\midrule
Component & Which step spends? & One tool dominates \\
Task type & Which workload spends? & A rare task type dominates \\
User / tenant & Who spends? & A small percentile dominates \\
Outcome & Was it spent well? & Failures cost more than successes \\
\bottomrule
\end{tabular}
\end{center}

The outcome dimension is the one teams skip and the one that pays. Cost per
\emph{successful} task is the metric that matters, and it behaves differently
from cost per request
\[
C_{\text{unit}} = \frac{C_{\text{total}}}{N_{\text{success}}}
= \frac{\bar{c}_{\text{req}}}{p_{\text{success}}}.
\]
A system with an 80\% success rate pays a $1.25\times$ multiplier on every
request; at 50\% it pays $2\times$. Reliability is a cost optimization, and
frequently a larger one than any efficiency measure, because failed runs consume
full price and produce nothing. Worse, failures skew expensive. A run that
fails at step 12 has paid for 12 steps.

Two costs are routinely omitted from these dashboards. Retries, which are
attributed to the successful attempt if they are attributed at all. And storage, including agent runs leave behind logs, context snapshots, checkpoints, and debug traces
that persist and accrue, a footprint that benchmarks and cost models have largely
ignored \cite{footprint2026}. For a fleet of long-running agents this is a real
line item rather than a rounding error.

\section{Worked Example: Cost-Optimized Research Agent}

A research agent answers analyst questions over an internal corpus. It handles
50{,}000 queries per month. The naive implementation routes everything to a
frontier model, retrieves ten documents per query, and replays the full history
across an average of four turns.

\textbf{Baseline.} Per query, namely 1{,}500-token system prompt, ten documents at 800
tokens, 200 tokens of question, and roughly 600 tokens of output per turn, with
history replayed each turn.

\begin{center}.
\begin{tabular}{lrrr}
\toprule
Component & Tokens/query & \$/query & \$/month \\
\midrule
System prompt (replayed $4\times$) & 6{,}000 & 0.015 & 750 \\
Retrieved documents ($10\times800$, replayed) & 20{,}000 & 0.050 & 2{,}500 \\
History accumulation & 4{,}800 & 0.012 & 600 \\
Output ($4 \times 600$) & 2{,}400 & 0.024 & 1{,}200 \\
\midrule
Total & 33{,}200 & 0.101 & 5{,}050 \\
\bottomrule
\end{tabular}
\end{center}

Now apply four interventions, measuring each against a held-out accuracy set
rather than assuming they are free.

\begin{center}.
\begin{tabular}{llrr}
\toprule
\# & Intervention & \$/month & Accuracy \\
\midrule
0 & Baseline & 5{,}050 & 84.0\% \\
1 & Prompt-cache the stable prefix ($h=0.9$) & 4{,}430 & 84.0\% \\
2 & Retrieve 4 documents, rerank instead of 10 & 3{,}180 & 84.6\% \\
3 & Cascade: small model first, escalate 25\% & 1{,}560 & 83.1\% \\
4 & Compress history; cap output at 400 tokens & 1{,}160 & 82.3\% \\
\bottomrule
\end{tabular}
\end{center}

Total reduction is 77\%, at a cost of 1.7 accuracy points.

Read the middle rows carefully, because they carry the lesson. Step~1 is free, such as identical outputs, lower bill, no reason not to do it. Step~2 \emph{improved}
accuracy while cutting cost by a third. Retrieving ten documents was introducing distractors, so the cheaper configuration was also the better one.
Steps~3 and~4 are genuine trades, spending accuracy for money.

So the frontier is not a smooth curve you slide along. It has free moves,
strictly better moves, and real trades, and they need to be identified separately.
A team that treats all four as ``cost optimization'' will negotiate about the
1.7-point drop without noticing that the first two steps cost nothing at all.

Whether the last two steps are acceptable depends on what a wrong answer costs.
For exploratory analyst queries where the user sees the citations and can check,
82.3\% at \$1{,}160 is an easy yes. For a system whose output feeds an automated
decision, the 1.7 points may be worth \$400 a month several times over, and
Chapter~\ref{ch:gates} gives the machinery for deciding rather than guessing.

\section{The Cost-Quality Frontier}

Plot every configuration you have measured with cost on one axis and quality on
the other. Most points are dominated. Some other configuration is cheaper and
better. The undominated points form the frontier, and the only interesting
question is where on it to sit.

Two properties make this harder than a standard Pareto analysis.

The frontier moves. Inference prices fall by something like an order of magnitude
per year for a fixed capability level \cite{llmflation2024, epochpricing2025}.
A configuration that is unaffordable today is routine within a year, so a
sufficiently complex optimization can be obsolete before it finishes paying for
itself. Ask how long an optimization takes to build, maintain, and validate
against the possibility of buying the same saving by waiting.

That argues for a specific bias. Prefer optimizations that are cheap to build and
architectural rather than clever. Prompt-cache-friendly prompt layout keeps
paying. A hand-tuned distillation pipeline for one model version does not, and it
becomes a liability the moment that model is deprecated. Build for the ability to
swap in a cheaper model without changing the architecture, and let deflation do
the work that deflation does well.

The second property is that quality is not scalar. The worked example collapsed
it to one accuracy number, which hid that step~2 improved grounding while
step~3's escalation errors clustered on exactly the hard queries analysts cared
about most. Chapter~\ref{ch:evaluation} takes this up in earnest.

\section{Fallacies and Pitfalls}

\textbf{Fallacy. Optimize first, measure later.} Intuition about where an agent
spends money is reliably wrong. The usual culprit is history replay, which is
invisible in the code because nobody wrote a line that says ``send the entire
conversation again.''

\textbf{Fallacy. Input tokens are the bill.} Output costs several times input per
token. A verbose response format can dominate a large retrieved context.

\textbf{Fallacy. Mean cost describes the system.} Cost distributions across agent
tasks are heavy-tailed. A p99 run that takes 40 steps instead of 4 can carry a
substantial share of the monthly total, and the mean will not show it. Budget
against tail behavior, and read Chapter~\ref{ch:evaluation} on why the mean is
the wrong summary for almost everything an agent does.

\textbf{Fallacy. Failures are cheap because they are incomplete.} Failed runs pay
full price for every step they took, and they take more steps than successes
because agents retry. Cost per success is the honest denominator.

\textbf{Fallacy. Cache hit rate measures cache value.} A semantic cache with an
excellent hit rate and a loose threshold is a machine for confidently returning
answers to questions nobody asked.

\textbf{Pitfall. Cost attribution without an outcome field.} Without knowing
whether a run succeeded, a dashboard can only rank components by spend, never by
waste.

\textbf{Pitfall. Optimizing latency and cost with one lever.} Batching improves
cost and worsens latency. Speculation improves latency and worsens cost. Only
caching improves both, which is why it comes first.

\section{Summary}

Cost in an agent system is an accumulation problem, and accumulation is invisible
until it is instrumented.

Token costs dominate, and multi-turn history replay makes them grow
super-linearly in the number of turns. Decompose the total into tokens, tools,
compute, and the storage that agent runs leave behind, because each responds to a
different fix. Divide by successes rather than requests, since reliability is a
cost lever and usually a bigger one than efficiency.

Latency follows the critical path rather than the total, so parallel work is free
in latency terms and never free in cost terms. Caching is the only intervention
that improves both, and it comes in three layers that exploit three different
redundancies. The rest trades quality or latency for money. Cascading, context management, output control, and early stopping all belong in that group, and those trades should be measured rather than assumed.

Finally, the whole frontier slides downward about an order of magnitude a year.
Optimizations that are architectural keep paying. Optimizations that are clever
and version-specific are usually a slower way of buying what deflation would have
given you for free.

Chapter~\ref{ch:budgets} takes the next step. Making the budget a state variable
the policy can see and act on, rather than a limit it discovers by hitting.

\section*{Bibliographic Notes}

The LLMflation analysis documenting 10$\times$ annual cost reduction appears in \cite{llmflation2024}. Epoch AI's analysis \cite{epochpricing2025} shows 9$\times$ to 900$\times$ annual price decline depending on capability tier, with median 200$\times$ in recent periods.

KV caching is standard in transformer inference; PagedAttention \cite{vllm2023} introduced efficient memory management now adopted across inference frameworks. Prompt caching implementations are documented by Anthropic, OpenAI, and Google in their API documentation.

BudgetMLAgent \cite{budgetmlagent2024} demonstrates 94\% cost reduction through model cascading. AgentDiet \cite{agentdiet2024} shows trajectory compression for agent cost reduction. The State of Agent Engineering survey \cite{stateofagent2025} documents that latency has become the second-largest challenge (20\% of respondents) while cost concerns have decreased.

Critical path analysis for distributed systems is classical; application to LLM agents adapts these concepts to the unique characteristics of autoregressive generation (sequential output) and multi-turn interactions.

\section*{Exercises}

\begin{enumerate}
\item \textbf{Cost Calculation}. An agent uses GPT-4o (\$2.50/M input, \$10/M output) with a 1,500-token system prompt. Each query averages 200 input tokens and 400 output tokens. The agent makes 3 LLM calls per task with full context replay. Calculate, including (a) cost per task, (b) cost for 50,000 tasks/month, (c) savings if prompt caching achieves 90\% hit rate on system prompt.

\item \textbf{Critical Path Analysis}. A multi-agent system has, namely Router (0.5s), two parallel workers (2s and 3s), Aggregator (1s), and final response (1.5s). Calculate, such as (a) critical path latency, (b) latency if workers were sequential, (c) latency reduction from parallelization.

\item \textbf{Model Cascade Design}. Design a cascading strategy for a customer service agent handling 100,000 queries/day. Given, including GPT-4o-mini (\$0.15/M in, \$0.60/M out, 80\% accuracy), GPT-4o (\$2.50/M in, \$10/M out, 95\% accuracy). Estimate, namely (a) cost with only GPT-4o, (b) cost with cascade (try mini first, escalate failures), (c) expected accuracy with cascade.

\item \textbf{Cache Break-Even}. A prompt caching system has, such as cache write cost 1.25$\times$ base, cache read cost 0.1$\times$ base, cache hit rate $h$, system prompt 3,000 tokens. Derive, including (a) formula for cost savings as function of $h$, (b) minimum $h$ for positive savings, (c) savings at $h = 0.8$.

\item \textbf{LLMflation Planning}. A system costs \$50,000/month today. Assuming 10$\times$ annual cost reduction for equivalent capability: (a) project costs for years 1, 2, 3, (b) calculate NPV of a \$100,000 optimization investment that reduces current costs by 50\%, (c) recommend whether to invest or wait.
\end{enumerate}

\section*{Bibliography}

\begin{thebibliography}{99}

\bibitem{agentdiet2024}
Yuan-An Xiao, Pengfei Gao, Chao Peng et al.
\newblock Reducing Cost of LLM Agents with Trajectory Reduction.
\newblock {\em arXiv:2509.23586}, 2025.

\bibitem{anthropicresearch2025}
Anthropic.
\newblock How We Built Our Multi-Agent Research System.
\newblock Engineering blog, 2025.

\bibitem{budgetmlagent2024}
A. Gandhi et al.
\newblock BudgetMLAgent: A Cost-Effective LLM Multi-Agent System for Automating Machine Learning Tasks.
\newblock {\em arXiv:2411.07464}, 2024.

\bibitem{epochpricing2025}
B. Cottier, B. Snodin, D. Owen, T. Adamczewski.
\newblock LLM Inference Prices Have Fallen Rapidly but Unequally Across Tasks.
\newblock Epoch AI, March 2025.

\bibitem{llmflation2024}
G. Fang, M. Casado.
\newblock Welcome to LLMflation: LLM Inference Cost is Going Down Fast.
\newblock Andreessen Horowitz, November 2024.

\bibitem{stateofagent2025}
LangChain.
\newblock State of Agent Engineering.
\newblock Survey report, December 2025.

\bibitem{vllm2023}
W. Kwon, et al.
\newblock Efficient Memory Management for Large Language Model Serving with PagedAttention.
\newblock {\em SOSP}, 2023.


\bibitem{abortedkv2026}
Guijia Zhang and Harry Yang.
``Aborted but Not Forgotten: KV-Cache Retention Breaks Rollback Consistency in Language Agents.''
arXiv:2608.15939, 2026.

\bibitem{budgettooluse2026}
Tengxiao Liu et al.
``Budget-Aware Tool Use Enables Effective Agent Scaling.''
arXiv:2511.17006, 2025. COLM 2026.

\bibitem{footprint2026}
Chenglin Yu et al.
``The Hidden Footprint: Making Storage a First-Class Metric for LLM Agent Evaluation.''
arXiv:2607.11149, 2026.

\bibitem{kvsurvey2026}
Jiantong Jiang et al.
``Towards Efficient Large Language Model Serving: A Survey on System-Aware KV Cache Optimization.''
arXiv:2607.08057, 2026. ACL 2026.

\bibitem{notation2026}
Lorenz Kutschka and Bernhard Geiger.
``Notation Matters: A Benchmark Study of Token-Optimized Formats in Agentic AI Systems.''
arXiv:2605.29676, 2026.

\bibitem{progrouter2026}
Songyuan Li et al.
``ProgRouter: Online Progress-Guided Orchestration for Multi-Agent LLM Workflows under Quality-Cost Tradeoffs.''
arXiv:2608.25992, 2026. EMNLP 2026.

\bibitem{scmas2026}
Di Zhao et al.
``SC-MAS: Constructing Cost-Efficient Multi-Agent Systems with Edge-Level Heterogeneous Collaboration.''
arXiv:2601.09434, 2026.

\bibitem{triaxialkv2026}
Hanzhang Shen et al.
``TriAxialKV: Toward Extreme Low-Precision KV-Cache Quantization for Agentic Inference Tasks.''
arXiv:2605.17170, 2026.

\bibitem{ucci2026}
Varun Kotte.
``UCCI: Calibrated Uncertainty for Cost-Optimal LLM Cascade Routing.''
arXiv:2605.18796, 2026.

\bibitem{kvreuse2026}
Sichu Liang et al.
``When KV Cache Reuse Fails in Multi-Agent Systems: Cross-Candidate Interaction is Crucial for LLM Judges.''
arXiv:2601.08343, 2026.

\bibitem{tcoanalysis2024}
Amit Sharma, Teodor-Dumitru Ene, Kishor Kunal et al.
``Assessing Economic Viability: A Comparative Analysis of Total Cost of Ownership
for Domain-Adapted Large Language Models versus State-of-the-art Counterparts.''
arXiv:2404.08850, 2024.

\bibitem{plancache2025}
Qizheng Zhang, Michael Wornow, and Gerry Wan.
``Agentic Plan Caching: Test-Time Memory for Fast and Cost-Efficient LLM Agents.''
arXiv:2506.14852, 2025.

\bibitem{backpressure2026}
Zuyuan Zhang et al.
``Augmented Backpressure for Decentralized Management of Agentic Networks.''
arXiv:2608.00914, 2026.

\bibitem{cachecollision2026}
Zhixiang Zhang et al.
``From Similarity to Vulnerability: Key Collision Attack on LLM Semantic Caching.''
arXiv:2601.23088, 2026. ICML 2026.

\bibitem{saga2026}
Dongxin Guo et al.
``SAGA: Workflow-Atomic Scheduling for AI Agent Inference on GPU Clusters.''
arXiv:2605.00528, 2026.
\end{thebibliography}
% ============================================================
% ============================================================
% Chapter 2: Cost Models and Critical Paths
% ============================================================
% Chapter: Cost Models and Critical Paths, Part II
% ================================================

